\subsection*{一步一步建立直观}

\paragraph{1. 从乘法结构寻找圆上的指数}

实指数把相加的参数变成相乘的倍率；圆周旋转也有同样结构，角度相加等于旋转连续执行。因此用 $e^{i\theta}$ 表示单位圆位置，与指数的组合规律相容。右图让你先把这个结构看成动作，再与正弦余弦坐标对应。



\paragraph{2. 特殊角度只是一般规律的取样}

在半圈处读到负一，在整圈处回到一。这些值帮助检查符号，却不能独自证明所有角度下的欧拉公式。可以通过幂级数、微分方程等方法把一般规律建立起来；本节正文采用的论证应与它依赖的定义相互配套。



\paragraph{3. 指数的实部负责模长}

第三步让指数同时具有实部和虚部，轨迹离开单位圆。实部带来指数缩放，虚部带来旋转。理解了这种分工，就不会把复指数的“周期性”错误地理解成所有复方向都周期；周期出现在纯虚方向的整圈平移中。


